\chapter{文献综述与理论基础}
\label{chap:litreview}

\section{数字经济测度研究综述}

数字经济测度大致沿两条路径展开：一是基于增加值核算的规模测算，二是基于综合指标体系的水平评价\citep{xugc2023digital}。前者口径清晰但对统计数据要求较高；后者灵活性强、便于跨地区横向比较，但权重确定方法对结果影响较大\citep{liwei2022region}。本文采用综合指标体系路径，并以客观赋权缓解主观随意性。

\section{区域差异的测度与分解}

刻画区域差异的常用工具包括基尼系数、变异系数与泰尔指数。其中泰尔指数因满足可加分解性而被广泛采用：总差异可分解为组内差异与组间差异之和。设地区 $i$ 的数字经济发展水平占总体的比重为 $s_i$，地区个数为 $n$，则泰尔指数定义为

\begin{equation}
    T = \sum_{i=1}^{n} s_i \ln\!\frac{s_i}{1/n}.
    \label{eq:theil}
\end{equation}

式~\eqref{eq:theil} 取值越大，表示区域间发展越不均衡。当样本被划分为若干组时，$T$ 可进一步分解为组间项与组内项之和，从而识别差异的主要来源。

\section{理论基础}

数字经济的区域非均衡可由要素集聚与空间溢出理论加以解释：数字基础设施与人才在中心地区集聚，形成规模报酬递增，进而拉大区域差距\citep{chen2019cluster}。这一理论视角为后文的差异分解提供了解释框架。
